%%% FP2001.tex 05.10.2000 - Vesa Kolhinen <vesa.kolhinen@phys.jyu.fi>
%%% Original version:
%%% FP2000.tex 12.10.1999 - Eero Noponen <eero.noponen@hut.fi>
%%%
%%%
%%%
%%% Sample LaTeX file for the abstracts for the Proceedings of the
%%% Annual Meeting of the Finnish Physical Society.

\documentclass[12pt]{article}
\usepackage{fp2009}
\usepackage[T1]{fontenc}
\usepackage[latin1]{inputenc}
\usepackage{hyperref}

%%% Remove the following line if Times font is unavailable.
\usepackage{times}

%%% The following line is required if the document contains figures.
\usepackage{graphicx}

%%% The following line is for LaTeX 2.09.
%\documentstyle[FP2001,12pt]{article}

\begin{document}

%%%
%%% The title of the abstract is given here as \title{TITLE TEXT}.
%%%
\title{INCL/ABLA INTRA-NUCLEAR CASCADE MODEL UPDATES}

%%%
%%% The list of authors; underline the name of the author who is responsible
%%% for the presentation of the paper.
%%%
A. Heikkinen \underline{P. Kaitaniemi}

%%%
%%% The address.
%%%
P.O.Box 64 (Gustaf H\"allstr\"omin katu 2), FIN-00014 University of Helsinki, Finland\\
email: pekka.kaitaniemi@gmail.com

%%%
%%% Body text begins here.
%%%

We present a new Geant4 physics list prepared for nuclear physics applications
in the domain dominated by spallation.
We discuss new Geant4 models based on the translation of 
INCL intra-nuclear cascade and ABLA de-excitation codes in C++   
and used in the physic list.
The INCL model is well established for targets heavier than Aluminium
and projectile energies from $\sim$ 150 MeV up to 2.5 GeV $\sim$ 3 GeV. 
Validity of the Geant4 physics list is demonstrated from the perspective of accelerator driven systems
and EURISOL project, especially with the neutron double differential cross sections and residual
nuclei production.
Foreseen improvements of the physics models for the treatment of light targets (Carbon - Oxygen)
and light ion beams (up to Carbon) are discussed.
%Improvements of the physics models are foreseen for the treatment of light targets (Carbon - Oxygen)
%and light ion beams (up to Carbon) and will be discussed.
An example application utilizing the physics list is introduced.
%Outside this domain (light nuclei and energies down to 20-30 MeV)
%the model can be tried but the capabilities have yet to be established.
%such as spallation reactions in Accelerator Driven Systems.

%%% The following minipage structures are needed to position the figure
%%% beside the text paragraph.
%%% For a standard figure environment, see the example below.

\begin{minipage}{9cm}
Figures may be included in the submitted files. Here we show how a
figure can be included in a \LaTeX\ file.
\end{minipage}
\hfill
\begin{minipage}{5cm}
\includegraphics[width=4cm]{logo.eps}
\end{minipage}

%%%
%%% A more typical Figure.
%%%
%\begin{figure}[!h]
%\centering\includegraphics[width=4cm]{logo.eps}
%\caption{Mallikuva.}
%\end{figure}

%%%
%%% The bibliography: the references are listed here.
%%%
\begin{thebibliography}{9}

\bibitem{ref1}
A. Fyysikko and B. Fyysikko,
\href{http://link.aps.org/abstract/PR/v108/p1175}{Fysiikan Lehti 108
(1957) 1175}.
\end{thebibliography}

\end{document}
